
In this project, it was studied in depth a two-dimensional discrete-time dynamical system given by two logistic maps. The aim was to investigate how the long-term behaviour of the system changes for different parameter values, using fixed points, periodic points, phase portraits and a bifurcation diagram. Which was completed sucessfully.

It is shown a typical progression of nonlinear dynamical systems for this model, which means that for smaller parameter values, trajectories may converge to a single attracting fixed point. The theoretical fixed points were calculated and match the posterior obtained results with the exception of the symmetric case $r_1=r_2=3.1$, where $x=y=1-\frac{1}{3.1}=0.67741935$ exists but due to not being attractive it does not show up in the numerical simulations. This system, therefore, converges to a stable period-2 orbit, from what we can conclude that the existence of a fixed point alone is not enough to describe the observed long-term behaviour, it's then clear that its stability is essential. 

Additionally, it was observed that the symmetry between the two populations is broken when increasing $r_2$, while keeping $r_1=3.1$. The systems exhibits a period-2 behavior which then turns into an asymmetric orbit. On the other hand, for larger values of r such as $r_2=3.8$, the bifurcation cascade and the computation of periodic points both indicate a stable period-4 orbit. With even greater values, such as $r_2=3.95$, the system no longer converges to a small number of points, but instead produces a point cloud, suggesting more complex and chaotic dynamics.

This is all to show that different numerical approaches give a consistent picture. All the numerical calculations and simulations act accordingly. This leads to the conclusion that deterministic systems can easily become more complex, generating seemingly random behaviors which can actually be predicted in a universal way.

